\section{What sets the band}\label{sec:costs}

\Cref{prop:cs} says that trading costs set the edges of the band. This section tests that prediction in two ways: over
time, using the three increases in the tax on futures, and in levels, by adding up what an arbitrage position actually
costs. It then asks whether the band is wider for stocks that are harder to trade.

\subsection{Tax increases (pre-specified test H3)}
The STT on the sale of futures rose on 3 April 2023 (the first trading day after the 1 April change), 1 October 2024 and
1 April 2026, by 0.25, 0.75 and 3.0\bp\ of the futures sale. Each increase raises the cost of cash-and-carry, so the
upper edge of the band should rise, and by more for larger increases. Because \cref{prop:inv} implies that fixed
thresholds measure an average edge rather than the cost, the plan tests the \emph{change} in the upper part of each
stock's basis distribution. For stock $i$ and event $e$, let $q^{80}_{ie,\text{pre}}$ and $q^{80}_{ie,\text{post}}$ be
the 80th percentiles of its basis in the 40 trading days before and the 40 from the event. The test statistic is the
average change
\begin{equation}\label{eq:h3}
\overline{\Delta q}^{80}=\frac{1}{N}\sum_{i,e}\big(q^{80}_{ie,\text{post}}-q^{80}_{ie,\text{pre}}\big),
\end{equation}
over the 2024 and 2026 events, with standard errors clustered by stock. The 2023 event is reported but not used as
evidence, because the taxation of debt funds changed on the same day, which may have moved arbitrage capital at the
same time.

The upper percentile rose after both later increases, by 6.4\bp\ in 2024 and 7.4\bp\ in 2026, and fell slightly in
2023 (\cref{tab:h3}). The pooled change is $+6.96$\bp\ (standard error 0.61), and the three events are ordered as the
cost changes are. On the test fixed in advance, H3 is supported, and it remains significant after the Holm correction for
the four primary hypotheses.

\begin{table}[t]
\caption{The upper part of the basis rose after the 2024 and 2026 tax increases}\label{tab:h3}
\exsub{Mean change in each stock's percentiles of the basis, 40 trading days after the event against 40 before (bp)}
\small\begin{tabular}{@{}>{\bfseries\color{accent}}lccccc@{}}
\toprule
\thd{Event} & \thd{STT change} & \thd{Stocks} & \thd{80th percentile} & \thd{20th} & \thd{Median} \\
\midrule
3 Apr 2023$^{a}$ & $+0.25$ & 181 & $-2.12$ (0.87) & $+0.24$ & $+0.27$ \\
1 Oct 2024 & $+0.75$ & 172 & $+6.42$ (0.70) & $+4.55$ & $+5.39$ \\
1 Apr 2026 & $+3.0$ & 192 & $+7.44$ (0.98) & $+2.10$ & $+6.30$ \\
\bottomrule
\end{tabular}

\exsource{Standard errors in brackets. STT change in bp of the futures sale. $^{a}$ Same day as the change in the taxation of debt
funds; reported, not used as evidence. Source: author's calculations.}
\end{table}

Two features of the result call for caution. First, the rises are larger than the cost changes: an increase of 0.75\bp\
in tax went with a 6.4\bp\ rise in the upper percentile. Second, the median and the 20th percentile also rose in 2024 and
2026, so the whole distribution shifted rather than just its upper edge. Such shifts can come from anything that affects
all stocks at once, and a standard error clustered only by stock does not allow for them, because every stock shares the
same event date (Kolari and Pynn\"onen, 2010). The plan therefore also specified a placebo check: the same statistic at
30 dates away from the true events. Across these dates it has a mean of $-1.0$\bp\ and a standard deviation of 7.4\bp,
and 10\% of placebo dates show a rise at least as large as in 2024, 7\% one as large as in 2026 (\cref{fig:h3placebo}).
Judged against what happens on ordinary dates, the rises after the tax increases are large but not exceptional. My
reading is that costs push the band in the predicted direction but do not set its edges as tightly as the model
implies.\footnote{The same statistic at the three dates of SEBI's 2024 to 2025 reforms of index derivatives, which did not
change the cost of single-stock arbitrage, gives $-4.1$, $-0.3$ and $+1.9$\bp, all within the placebo range.}

\begin{figure}[t]
\caption{Against placebo dates, the rise after the tax increases is large but not exceptional}\label{fig:h3placebo}
\exsub{Mean change in stocks' 80th percentile of the basis at the three STT dates and at 30 placebo dates (bp)}
\centering\includegraphics[width=\linewidth]{fig_h3_placebo.pdf}
\exsource{Placebo dates every 80 trading days, away from the true events. Source: author's calculations.}
\end{figure}

\subsection{What an arbitrage position costs}
A second way to compare the band with costs is to add up the cost of an actual position. An overnight cash-and-carry
position held to delivery pays STT of 0.1\% on the purchase of the stock and again on delivery at expiry, the STT on the
futures sale, stamp duty, exchange charges, the SEBI fee, goods and services tax on fees, and brokerage. With an
institutional brokerage of 0.01\% per leg, the full round trip costs 27.1, 27.3, 28.0 and 31.0\bp\ in the four tax
periods (\cref{fig:costband}), or 25 to 38\bp\ for brokerage between 0.005\% and 0.03\% per leg. Rolling an existing
position for one more month is far cheaper, 4 to 8\bp, because the delivery taxes are paid only once.

\begin{figure}[t]
\caption{The full cost of a new position lies above the estimated upper edge}\label{fig:costband}
\exsub{Round trip of an overnight cash-and-carry position held to delivery, by STT period (bp of notional); triangles: estimated upper threshold in levels}
\centering\includegraphics[width=\linewidth]{fig_cost_band.pdf}
\exsource{Institutional brokerage of 0.01\% per leg; GST of 18\% on fees and brokerage. Source: NSE circulars on STT and transaction charges; author's calculations.}
\end{figure}

The estimated upper threshold of the basis, measured in levels rather than relative to each stock's median, lies between
$+22$ and $+28$\bp\ in the four periods, below the full cost of a new position, which falls between the 90th and 95th
percentiles of the basis. This is what \cref{prop:inv} implies when arbitrageurs already hold positions. New capital
enters only when the gap is large enough to pay the full round trip, while funds that keep their capital in place and
roll it from month to month can lean against much smaller gaps. The edge measured in the data is a mixture of the two.

\subsection{Frictions across stocks (H1)}
The band should also be wider for stocks in which arbitrage is harder. For each F\&O stock and quarter I build a
friction index from five measures taken in the previous quarter: the bid-ask spread estimated from daily high and low
prices (Corwin and Schultz, 2012), illiquidity (Amihud, 2002), idiosyncratic volatility, utilisation of the position
limit, and the share of days with no stock lending. The index is the first principal component of these measures,
standardised within each quarter. I then regress the width of each stock's band in the quarter (the 80th minus the 20th
percentile of its basis on non-ban days) on the index, with quarter fixed effects and standard errors clustered by stock
and quarter.

A stock with an index one standard deviation higher has a band 3.5\bp\ wider (standard error 1.1, one-sided $p=0.002$;
the mean width is 40.8\bp) and a mean absolute basis 2.5\bp\ larger (\cref{fig:h1bins}). Most of the difference is
between the least frictional stocks and the rest. The relation holds across stocks, not within a stock over time: with
stock fixed effects it disappears, so it reflects persistent differences between stocks rather than changes in a
stock's frictions. The index also does not predict how fast the basis returns to the band, which is consistent with the
finding above that daily reversion is mostly noise. One component enters with the opposite sign: stocks with fewer
lending days have narrower bands, which suggests that lending activity measures demand to borrow more than the ease of
borrowing.

\begin{figure}[t]
\caption{Bands are wider in stocks that are harder to trade}\label{fig:h1bins}
\exsub{Band width and mean absolute basis by decile of the friction index of the previous quarter, F\&O stocks, 2015 to 2025 (quarter means removed)}
\centering\includegraphics[width=\linewidth]{fig_h1_bins.pdf}
\exsource{Index: first principal component of spread, illiquidity, idiosyncratic volatility, utilisation of the position limit and lack of stock lending. 7,398 stock-quarters, 347 stocks. Source: author's calculations.}
\end{figure}
